% 题干转写：完成
% 解析转写：完成
% 图形核对：完成
\section{第36届复赛第5题：阶跃型光纤耦合与全反射}
\subsection{题干}
阶跃型光纤内芯直径为 $50\,\mu\mathrm{m}$、折射率 $n_1=1.46$；玻璃外包层外径为 $125\,\mu\mathrm{m}$、折射率 $n_2=1.45$。现用直径 $D=3.00\,\mathrm{mm}$、折射率 $n=1.50$ 的微球透镜将氦氖激光器输出的圆柱形平行光束耦合入光纤，空气折射率 $n_0=1.00$，激光束轴、透镜球心和光纤轴重合，如图\ref{fig:cpho36q05}。求：（1）光纤端面的最大入射角；（2）平行光刚好会聚在光纤端面时，微球后表面顶点 $O$ 与端面的距离；（3）满足长距离传输条件的平行入射光束最大直径。
\begin{figure}[htbp]\centering
\begin{tikzpicture}[scale=0.91,>=Stealth,font=\small]
  \draw[densely dashed,gray] (-4.35,0) -- (4.75,0);
  \draw[->,very thick] (-4.4,1.08) -- (-2.65,1.08);
  \draw[->,very thick] (-4.4,-1.08) -- (-2.65,-1.08);
  \node[left] at (-4.4,0) {平行激光束};
  \draw[fill=cyan!10,thick] (-1.85,0) circle (1.72);
  \node[above] at (-1.85,1.88) {微球透镜};
  \node at (-1.85,-1.15) {$n=1.50$};
  \coordinate (A) at (-2.95,1.08);
  \coordinate (B) at (-2.95,-1.08);
  \coordinate (F) at (1.02,0);
  \draw[thick] (-4.4,1.08) -- (A) -- (F);
  \draw[thick] (-4.4,-1.08) -- (B) -- (F);
  \fill (F) circle (1.2pt) node[below] {$F$};
  \fill (-0.13,0) circle (1.2pt) node[below left] {$O$};
  \draw[thick] (1.55,-0.92) rectangle (4.55,0.92);
  \draw[fill=gray!20,thick] (1.55,-0.30) rectangle (4.55,0.30);
  \node at (3.05,0) {$n_1=1.46$};
  \node at (3.05,0.61) {$n_2=1.45$};
  \node at (3.05,-0.61) {$n_2=1.45$};
  \node[above] at (3.82,1.05) {光纤};
  \node[above] at (2.25,1.05) {芯径 $50\,\mu\mathrm{m}$};
  \draw[<->] (4.92,-0.92) -- (4.92,0.92) node[midway,right] {$125\,\mu\mathrm{m}$};
  \draw[<->] (-0.13,-1.92) -- (1.55,-1.92) node[midway,below] {$OF=0.75\,\mathrm{mm}$};
\end{tikzpicture}

\caption{第36届复赛第5题微球透镜耦合阶跃型光纤的 TikZ 重绘。}\label{fig:cpho36q05}
\end{figure}

\subsection{参考解答整理}
\paragraph{（1）端面最大入射角。} 设端面最大入射角为 $\theta$，进入纤芯后的折射角为 $\alpha$。由端面折射和芯--包层界面的全反射临界条件，
\[
n_0\sin\theta=n_1\sin\alpha,
\qquad
n_1\cos\alpha=n_2.
\]
相应的端面折射与芯--包层全反射临界光线见图\ref{fig:cpho36q05solution}左图。
消去 $\alpha$，得
\[
\sin\theta=\sqrt{\left(\frac{n_1}{n_0}\right)^2-\left(\frac{n_2}{n_0}\right)^2}
=\sqrt{1.46^2-1.45^2}\approx0.17.
\]
因此
\[
\boxed{\theta_{\max}=\arcsin(0.17)\approx9.8^\circ.}
\]

\paragraph{（2）微球后顶点到端面的距离。} 对微球前表面，平行光经球面折射后的一次像距为
\[
\frac{n_0}{o_1}+\frac{n}{i_1}=\frac{n-n_0}{r},
\qquad o_1=\infty,\quad r=\frac D2,
\]
故
\[
i_1=\frac{nD}{2(n-n_0)}=4.5\,\mathrm{mm}.
\]
相对于后表面顶点 $O$，该像为虚物，
\[
o_2=i_1-D=\frac{(n-2n_0)D}{2(n-n_0)}=-1.50\,\mathrm{mm},
\qquad r_2=-\frac D2.
\]
再对后表面应用球面折射公式，得到二次像距
\[
\boxed{i_2=\frac{(2n_0-n)D}{4(n-n_0)}=0.75\,\mathrm{mm}.}
\]
因此 $O$ 到光纤端面的距离为 $0.75\,\mathrm{mm}$。几何近轴解法同样给出 $OF=D/4=0.75\,\mathrm{mm}$。
\begin{figure}[htbp]\centering
\begin{tikzpicture}[x=.7cm,y=.7cm,>=Stealth,font=\scriptsize]
  \begin{scope}
    \draw[densely dashed] (-4,0)--(3.7,0);
    \draw[thick] (0,-1)--(3.1,-1)--(3.1,1)--(0,1)--cycle;
    \draw[thick] (0,.25)--(3.1,.25);
    \draw[thick] (0,-.25)--(3.1,-.25);
    \node at (1.55,.62) {$n_2$}; \node at (1.55,0) {$n_1$}; \node at (1.55,-.62) {$n_2$};
    \draw[->,thick] (-2.2,.75)--(0,.25);
    \draw[->,thick] (0,.25)--(1.1,-.25);
    \draw[->,thick] (1.1,-.25)--(2.65,.25);
    \draw (0,.04) arc(0:22:.42) node[right]{$\theta$};
    \draw (0.35,.24) arc(0:-24:.42) node[below right]{$\alpha$};
    \node[below] at (1.55,-1.2) {(a) 端面折射与全反射临界光线};
  \end{scope}
  \begin{scope}[xshift=10cm]
    \draw[densely dashed] (-4,0)--(5.2,0);
    \draw[thick] (0,0) circle (1.65);
    \coordinate (Oc) at (0,0);
    \fill (Oc) circle(1.3pt) node[below]{$O'$};
    \fill (1.65,0) circle(1.3pt) node[below]{$O$};
    \coordinate (A) at (-1.05,.95); \coordinate (B) at (1.42,.82); \coordinate (F) at (3.0,0); \coordinate (C) at (4.25,0);
    \fill (A) circle(1.3pt) node[above]{$A$}; \fill (B) circle(1.3pt) node[above]{$B$};
    \draw[->,thick] (-3.8,.95)--(A);
    \draw[thick] (A)--(B)--(F);
    \draw[dashed] (A)--(C); \draw[dashed] (B)--(C);
    \draw[dashed] (Oc)--(A) (Oc)--(B);
    \fill (F) circle(1.3pt) node[below]{$F$}; \fill (C) circle(1.3pt) node[below]{$C$};
    \node at (0,-1.05) {$n$};
    \draw (A) ++(-.32,0) arc(0:-35:.38) node[right]{$\alpha$};
    \draw (A) ++(.27,-.04) arc(-8:20:.44) node[above right]{$\gamma$};
    \draw (B) ++(.33,-.02) arc(-15:16:.4) node[right]{$\alpha$};
    \node[below] at (.55,-2) {(b) 微球折射的近轴几何};
  \end{scope}
\end{tikzpicture}

\caption{第36届复赛第5题端面临界光线与微球两次折射的几何辅助图。}\label{fig:cpho36q05solution}
\end{figure}

\paragraph{（3）可耦合的最大光束直径。} 微球的等效焦距为
\[
f_e=\frac{nD}{4(n-n_0)}=2.25\,\mathrm{mm}.
\]
当边缘光线在光纤端面的入射角恰为 $\theta_{\max}$ 时，利用近轴几何关系（忽略 $O[(d/D)^3]$ 项）有
\[
d_{\max}=\frac{nD\sin\theta_{\max}}{2(n-n_0)}.
\]
代入参数，
\[
\boxed{d_{\max}=0.76\,\mathrm{mm}.}
\]
故平行入射激光束的直径不能超过 $0.76\,\mathrm{mm}$。

\subsection{转写依据}
题干、答案公式和图形分别依据 \texttt{CPHO36-Q05.md}、\texttt{CPHO36-Q05\_formula\_check\_v1.csv} 与题图核对入口。评分分布为第（1）问8分、第（2）问20分、第（3）问12分。
